% 分类目录：dl
\documentclass[12pt, a4paper]{article}
\usepackage[margin=2.5cm]{geometry}
\usepackage{ctex}
\usepackage{amsmath, amssymb}
\usepackage{listings}
\usepackage{xcolor}
\usepackage{tabularx}
\usepackage{hyperref}

\lstset{
    language=Python,
    basicstyle=\ttfamily\small,
    keywordstyle=\color{blue},
    commentstyle=\color{green!50!black},
    stringstyle=\color{red},
    showstringspaces=false,
    breaklines=true,
    numbers=left,
    numberstyle=\tiny\color{gray},
    frame=single
}

\title{知识点精讲：深度学习损失函数 (Loss Functions) 详解}
\author{}
\date{}

\begin{document}
\maketitle

\section{损失函数的统一数学框架}
从概率论视角看，损失函数通常是**极大似然估计 (MLE)** 的负对数似然形式。
\begin{itemize}
    \item \textbf{高斯分布假设 $\to$ MSE}：若假设模型噪声服从正态分布 $P(y|x) \sim N(f(x), \sigma^2)$，则最小化负对数似然等价于最小化均方误差。
    \item \textbf{伯努利分布假设 $\to$ 二元交叉熵}：若假设 $P(y|x)$ 服从二项分布，则 MLE 直接推导出 BCE 公式。
\end{itemize}

\section{信息论视角：从熵到交叉熵}
\begin{itemize}
    \item \textbf{自信息与熵}：$H(P) = -\sum P(x) \log P(x)$，衡量分布的不确定性。
    \item \textbf{KL 散度 (相对熵)}：$D_{KL}(P||Q) = H(P, Q) - H(P)$，衡量分布 $Q$ 对真实分布 $P$ 的偏离程度。
    \item \textbf{交叉熵}：$H(P, Q) = -\sum P(x) \log Q(x)$。在机器学习中，由于真实分布 $P$ 是固定的，最小化 KL 散度等价于最小化交叉熵。
\end{itemize}

\section{常用损失函数详解}

\subsection{回归任务 (Regression)}
\begin{enumerate}
    \item \textbf{均方误差 (MSE)}：对异常值极度敏感（平方惩罚）。
    \item \textbf{平均绝对误差 (MAE)}：鲁棒性更好，但在 0 点不可导。
    \item \textbf{Huber Loss}：误差小时平方，误差大时线性。具备 MSE 的平滑收敛与 MAE 的离群点抗性。
    \item \textbf{Log-Cosh Loss}：$\ln(\cosh(y - \hat{y}))$。二阶可导的 Huber 近似，比 Huber 更平滑。
\end{enumerate}

\subsection{分类任务 (Classification)}
\begin{enumerate}
    \item \textbf{交叉熵 (CE)}：分类标准，具有加速梯度传导的特性。
    \item \textbf{Focal Loss}：旨在解决类别不平衡和难易样本失配。
    $$ L_{focal} = -(1 - p_t)^\gamma \log(p_t) $$
    通过 $(1-p_t)^\gamma$ 衰减易分类样本的分值，使模型专注于“难点”。
\end{enumerate}

\subsection{计算机视觉与度量学习 (CV \& Metric)}
\begin{enumerate}
    \item \textbf{Dice Loss}：专门应对高度类别不平衡（如医学图像分割），衡量两个区域的重叠度。
    \item \textbf{IOU 家族}：从 IOU 演进到 \textbf{GIOU} (非重叠区补偿)、\textbf{DIOU} (中心距惩罚) 到 \textbf{CIOU} (宽高比一致性)，是目标检测回归精度的保障。
    \item \textbf{Triplet Loss}：三元组损失 $\max(0, d(a,p) - d(a,n) + \text{margin})$，用于人脸识别等嵌入式空间学习。
\end{enumerate}

\section{数学推导：Softmax 加交叉熵的导数}
设 $z_i$ 为网络输出的 Logits，$p_i = \text{softmax}(z_i)$，损失函数 $L = -\sum y_j \log p_j$。
对于第 $i$ 个维度的求导：
$$ \frac{\partial L}{\partial z_i} = p_i - y_i $$
其中 $y_i$ 为真值标签（One-hot）。此公式展示了反向传播的优雅：梯度即为**预测值与真实值的差**。

\section{全维度对比概览表}
\begin{table}[h]
\centering
\small
\begin{tabularx}{\textwidth}{|l|c|X|l|}
\hline
\textbf{损失函数} & \textbf{鲁棒性} & \textbf{梯度行为} & \textbf{核心应用} \\
\hline
MSE & 差 & 随误差减小而衰减 (有利收敛) & 线性回归、坐标预测 \\
\hline
MAE & 好 & 梯度恒定 (小误差下易震荡) & 鲁棒回归 \\
\hline
Focal Loss & --- & 降低易分类样本权重 & 类别极度不平衡 \\
\hline
Dice Loss & --- & 区域相关性驱动 & 语义分割 \\
\hline
\end{tabularx}
\end{table}

\section{标签平滑 (Label Smoothing)}
将目标值修改为 $\hat{y} = (1 - \epsilon)y + \epsilon/K$，防止模型在训练中过于“自信”，起到正则化作用，减少过拟合。

\section{关键代码片段 (PyTorch)}
\begin{lstlisting}
import torch
import torch.nn.functional as F

# 1. 具备数值稳定性的 BCE
loss_bce = F.binary_cross_entropy_with_logits(logits, target)

# 2. 自定义 Focal Loss 简易逻辑
def focal_loss(logits, targets, alpha=0.25, gamma=2.0):
    ce_loss = F.binary_cross_entropy_with_logits(logits, targets, reduction='none')
    pt = torch.exp(-ce_loss) 
    focal_loss = alpha * (1 - pt)**gamma * ce_loss
    return focal_loss.mean()
\end{lstlisting}

\section{参考文献}
\begin{enumerate}
    \item Lin, T. Y., et al. (2017). \textit{Focal Loss for Dense Object Detection}.
    \item Goodfellow, I., et al. (2016). \textit{Deep Learning} (Chapter 6.2).
    \item Sudre, C. H., et al. (2017). \textit{Generalised Dice Overlap as a Deep Learning Loss Function for Highly Unbalanced Segmentations}.
\end{enumerate}

\end{document}
